\chapter[Writing, Memory, and Power]{How Writing Changed Memory and Power}
\section{A Four-Thousand-Year-Old Bad Review}
First pick up an object. It is not beside you but in a display case at London's British Museum: a palm-sized clay tablet, dark brown, rounded at the corners, densely impressed with little wedge-shaped marks, as though a bird had walked back and forth across it.

Made about 3,750 years ago, it was excavated at ancient Ur in present-day Iraq. Its script is cuneiform, its language Akkadian. You might expect something so ancient and solemnly inscribed to be a poem, prayer, or law code. It is none of these. It is a complaint letter.

The writer is Nanni; the recipient a copper merchant, Ea-nasir. In essence Nanni complains: your copper ingots are inferior; my servant made repeated trips and you treated him rudely; you took payment but gave the good copper to others. From now on, I will send my own inspector and reject anything substandard.

In other words, a four-thousand-year-old ``bad review.'' Its vivid, righteously indignant tone has made it an internet celebrity in recent years, nicknamed ``the oldest surviving consumer complaint.''

Pause to appreciate the oddity. Writing, often called humanity's greatest invention and honored by civilizations as a sacred gift, spent much of its infancy not recording great minds like Homer's, but noting who delivered how much barley, how many sheep a warehouse held, and ``Your copper is terrible, and I am angry.''

This chapter explores the whole story behind that object: where writing came from, how it changed remembering, and how it quietly rewrote ``who decides.'' Memory and power, you will find, have been tangled together since writing's birth. Today they remain tangled in the phone you hold.

\section{An Apprentice with a Reed Pen}
Come into a scene.

Time: an early morning some four thousand years ago. Place: a large temple compound in Mesopotamia, on the plain between the Tigris and Euphrates. The air smells of river water, dried reeds, and the sheep enclosure.

Dawn has barely arrived, yet a queue already waits at the warehouse. Farmers carry baskets of barley. The first announces his name and basket count loudly. Beside him sits a young scribe with a shaved head, silent and bent over a half-dry tablet, holding a sharpened reed. He presses its angled face into clay, lifts, turns: a small wedge-shaped pit. Combinations make names; rows make numbers. Barley enters storage, stroke by stroke fixed on the tablet. The farmer leans forward to look. Without raising his head, the young man waves him back.

The scribe was not born knowing this craft. As a teenager he entered a city school called a ``tablet house'' by the Sumerians. Countless excavated practice tablets let us reconstruct daily life there. Students copied a master's word lists character by character onto a tablet's left side, smoothing mistakes away and starting again. Failure to recite a list, poor posture, or speaking out of turn could earn a beating. One surviving student composition complains in detail of several beatings in a day. The lists' content is revealing: object names first, wood, reeds, livestock; then official posts and titles; finally legal sentences and bureaucratic formulas. Notice the order. The school taught not ``express yourself,'' but ``record this kingdom.''

Those who endured entered a different life. Literacy belonged to a tiny privileged class in these city-states. Scribes avoided fieldwork and heavy labor, wore clean clothes, worked in shaded corridors, and entered temples and palaces. More subtle was the power in their hands.

Imagine the barley farmer's feelings. He cannot read. This year's delivery, last year's, what he still owes: all are recorded on someone else's tablets. If the tablet says ten baskets remain due, then ten remain due. He cannot refute it because the memory belongs not to him but to the tablet and its keeper. ``That is what the accounts say'' can decide whether a family eats or starves that winter.

Remember the picture: one person writes, another obeys. Much of this chapter asks what happened in that moment.

\section{Accounts First, or Poetry?}
First lay out the conflict, without rushing to an answer.

When discussing writing, we tend to imagine its most glorious uses: poetry, scripture, history, love letters. Asked why humans invented it, you might vote ``to transmit knowledge and culture.'' That natural picture makes writing a noble mind's tool.

Archaeological evidence tells almost another story. Mesopotamia's earliest known systematic writing, tablets from Uruk around 3200 BCE, overwhelmingly consists of inventories and ration records: barley, beer, livestock, allocations to workers. Early Egyptian hieroglyphs, a pictorial sign system carved on temples and tombs, appear widely on royal objects and monuments identifying kings and ownership. China's earliest known mature script is late-Shang oracle-bone writing, inscribed on turtle shells and animal bones for royal divination: will it rain tomorrow, is war auspicious, will the queen's child be a boy or girl? In Mesoamerica, Maya glyphs were chiefly carved on stelae recording dates of royal accessions, campaigns, and rituals.

In other words, no ancient civilization began its writing history with ``writing poetry.'' The opening chapters are almost always accounts, kings, and gods.

Our real question therefore emerges in two layers.

First: how did writing arise? If not for literature, was it for bookkeeping or divine communication? These routes yield different understandings of its nature: administrative tool or sacred ritual.

Second, deeper: is writing a tool for memory, or memory's replacement? Moving memory from brains to clay, paper, and screens, does it free people or hollow them out? Remember the farmer. When memory belongs only to writers, ``remembering'' becomes a gate controlled by power.

Do not answer yet. First hear how ancient people themselves argued.

\section{A King's Warning, a Scribe's Pride}
Ancient arguments over writing's value began earlier and ran fiercer than we imagine.

\textbf{Voice one: writing will destroy memory.}

In Plato's Phaedrus, Socrates tells an Egyptian story. The god who invented writing presents it to King Thamus, boasting that it is a remedy for memory and wisdom. The king refuses. In essence, he says writing will create forgetfulness in learners' souls: depending on external marks, they will stop recalling inwardly. They will remember signs, not the things themselves. You offer students wisdom's appearance, not wisdom.

It is easy to laugh this off as an old-fashioned objection, ironically preserved in a written book. But state the king's reasons fully: they are substantial.

First, what you truly remember belongs to you. Knowledge alive in your head can be recalled, combined, and applied in new situations. Knowledge existing only on paper leaves you empty-handed without the paper. Second, writing can repeat but cannot answer. Ask a teacher what a sentence means and they can explain or argue. Ask a book repeatedly and it only repeats itself. Written words are therefore easily taken out of context, with the author absent and unable to defend them. Third, oral transmission is not a lazy backup but a highly skilled living art. In societies without writing, memory is intensively trained. Greek bards recited epics of tens of thousands of lines; West African griots, hereditary oral historians and singers, can sing a family's genealogy and deeds across dozens of generations for hours, adapting to listeners. For them, memory is not a warehouse but a house rebuilt daily. Calling them ``uncivilized'' is arrogant and false. We will return to that prejudice.

None of these three reasons is empty. You experience them daily: a number saved in contacts is soon forgotten; a formula copied into notes but not understood remains unusable in an exam.

\textbf{Voice two: writing lets the dead keep speaking.}

Opposite stand the young warehouse scribe's colleagues. Mesopotamian teaching materials include boastful maxims, roughly ``What sort of scribe cannot even speak Sumerian?''---like today's ``What sort of programmer cannot even use English?'' Their pride had real grounds. Writing lets someone dead for three millennia complain ``in person'' about copper. A marriage agreement or land deed remains effective after both signatories become dust. A command leaves a king's mouth, crosses hundreds of kilometers, and reaches a border fortress unchanged. Writing frees information from the speaker's body, so it need not vanish with death, forgetting, or a change of story. Oral memory can never do that in the same way.

\textbf{Voice three: what about the majority who cannot write?}

A third voice was almost never recorded, because its speakers could not write. For most ancient people, writing was not a tool they held but something that descended upon them. Their debts, taxes, and sentences were recorded; their grievances were not. Writing did not abolish unequal memory, merely changed its carrier from ``who remembers well'' to ``who can write.'' The literate gained something others lacked: officially recognized memory. The shift was so deep that we still feel its echoes. ``In black and white'' remains weighty evidence; ``your word proves nothing'' remains a familiar refusal to the weak.

All three voices are before us. Does writing liberate or replace memory? Give power wings or put it on a leash? Before answering, we need a tool to understand what writing itself is.

\section[Thinking Bridge: Sound or Meaning?]{Thinking Bridge: Does a Sign Record Sound or Meaning?}
The biggest obstacle is assuming writing simply means ``drawn speech.'' First mark a firm boundary between writing and pictures:

\textbf{Pictures record things; writing records language.} You can understand a bison in a cave painting by looking, but cannot read it word for word as a sentence. It does not specify ``bison'' rather than ``a horned animal.'' Writing, by contrast, must correspond to a sequence of words in a language, with broadly consistent readings by those who know it. That is why archaeologists cautiously call marks on some pottery, such as those from China's Jiahu site, ``signs'' rather than ``writing.'' They may carry meaning without evidence that they record language. This boundary keeps the history of writing from becoming muddled.

Now the portable tool. For any writing system, ask: \textbf{which level of language do its basic signs represent?} Imagine language as a three-story building: meaning, morphemes and words, at the top; syllables in the middle, such as mā, ``mother''; phonemes at the bottom, where mā divides into an initial sound and a rhyme. Writing systems attach signs to different floors.

Those attached to the meaning level are often called morphographic systems. Chinese characters are a leading example. But here lies a widespread misconception, an easily misjudged counterexample: people imagine all Chinese characters are little pictures, representing meaning but not sound. This is wrong. Pictographs such as \mbox{日}, ``sun,'' and \mbox{山}, ``mountain,'' are only a small fraction. Most characters combine meaning and sound hints. In \mbox{妈}, mā, ``mother,'' \mbox{女}, ``woman,'' signals a female association; \mbox{马}, mǎ, ``horse,'' suggests the sound ma. In \mbox{铜}, tóng, ``copper,'' the metal component \mbox{钅} suggests meaning, while \mbox{同}, tóng, suggests pronunciation. When learning characters you already use sound to guess unfamiliar ones. Chinese writing contains a pronunciation-hint system, simply less explicit than an alphabet's.

Systems attached to syllables are syllabaries: one sign per syllable. Japanese kana are typical, each sign in the standard chart representing a syllable. Another moving example comes from the Cherokee, an Indigenous North American people. Early in the nineteenth century, Sequoyah, who could read and write no script himself, spent more than a decade designing over eighty syllabic signs for his language, completed around 1821. The system was simple enough to learn within days. Cherokee people rapidly gained literacy in large numbers and founded their own newspaper. Notice the counterintuitive point: someone unable to read invented a writing system through his own ingenuity.

Systems attached to phonemes are alphabets: a sign roughly represents a minimal sound unit. Most alphabets used today trace to a shared distant ancestor, Phoenician writing on the eastern Mediterranean coast about three thousand years ago. Moving west to Greece, it gained vowel signs and later developed into the Latin alphabet. Moving east and south, it gave rise to Aramaic, then Hebrew and Arabic scripts, and Brahmi, ancestor of India's scripts. Devanagari, used for Hindi, and Thai and Tibetan scripts belong to this family. Many apparently unrelated scripts on your phone are long-separated relatives.

There are special cases too, such as Korean Hangul. It is a phonemic alphabet whose letters form syllable blocks instead of simply lining up, making it look something like Chinese characters. Created under King Sejong in 1443, it is a deliberately designed script. We will return to the storm it stirred.

With the three-story tool, another popular misconception falls easily: ``Writing evolved from meaning signs to syllables to alphabets, the highest stage.'' Popular in nineteenth-century Europe, this does not stand. Chinese's supposedly inferior system has served one of the longest-continuing civilizations and still has over a billion users. Japanese successfully combines Chinese characters and two syllabaries. Alphabets are not a destination, only one fit between script and language. Alphabetic writing suits languages with clearly distinct consonants and vowels; syllabaries can be economical for languages with limited syllable inventories, like Japanese. Scripts differ by fit, not evolutionary rank. As in biology, fish are not inferior birds; they fit different environments.

\section{``An Eye for an Eye'' on a Stone Pillar}
Now read a little primary material and see what writing looks like on law's territory.

In the Louvre in Paris stands a black basalt pillar over two meters tall, from Old Babylon around 1754 BCE. A relief at its top shows King Hammurabi before Shamash, god of the sun and justice, receiving authority. Below it, dense cuneiform records the famous Code of Hammurabi, with about 280 surviving provisions.

In its prologue Hammurabi says the gods commissioned his legislation, broadly to establish justice, destroy evil, and prevent the strong from oppressing the weak. The provisions are direct. Section 196 reads:

\begin{quote}
If a man destroys another man's eye, his eye shall be destroyed.
\end{quote}
This is the literal source of the famous ``eye for an eye'' principle. Related provisions address broken bones and knocked-out teeth.

Read the pillar through this chapter's lens, and at least three layers appear.

First, written law moves judgment from ``a powerful person's mouth'' to stone. Under oral judgment, two farmers' dispute depends on an elder's recollection of precedent, and power can knead memory. An inscribed code is, in principle, the same for everyone; changing a law means confronting a block of basalt. Here writing really puts a leash on power.

Second, the king still holds the leash's other end. Notice the relief's design: law is not negotiated by the people but given by a god to the king, then bestowed on subjects. And the text is Akkadian cuneiform. How many people in the kingdom could read it? For most illiterate subjects, the pillar was less an instruction manual than a monument silently declaring ``I define justice.'' Writing elevated power through the very gesture of making it public.

Third, most counterintuitively, few readers made the pillar safe. What threatened royal power was not the code itself but the eventual day when ordinary people could read it, quote it, and turn it back on the king. That day would wait until literacy ceased to be a minority privilege. Its arrival is a story for later.

\section{Why Writing Began: Three Explanations}
Separate four kinds of content. What happened is relatively clear from surviving evidence: early Mesopotamian texts are mostly accounts, Chinese examples divinations, Maya examples royal commemoration. Sources also tell us how people understood writing: Sumerians called it a divine gift; Shang kings treated inscriptions as messages to ancestors. How to judge it is yours to decide. The contest concerns why: what was writing originally for?

\textbf{Explanation one: bookkeeping (the economic account).}

Archaeologist Denise Schmandt-Besserat proposed an influential chain. Millennia before writing, Mesopotamian farmers kept accounts with small clay tokens: a cone for a sack of grain, a sphere for an animal. Tokens were enclosed in sealed clay envelopes for transport. To show the contents without opening them, people impressed tokens on the outside. Eventually they simply impressed or incised tablets; tokens retired and marks became signs, producing cuneiform. The account's strength is a complete evidential chain, excavated objects at each step, matching the predominance of early accounts. Its limitation is geography: it chiefly explains Mesopotamia. Shang kings inscribed bones to question spirits, not count warehouse stocks; Maya royal stelae were not ledgers. An explanation can govern one civilization without governing all.

\textbf{Explanation two: gods and kingship (the ritual account).}

This view notices something bookkeeping passes over: sacred halos surround early writing in several civilizations. China's first mature examples are divinations; Egyptian hieroglyphs appear widely on temples, tombs, and royal objects, and Egyptians called writing ``the gods' words.'' Maya inscriptions are almost entirely royal affairs: accessions, campaigns, rituals. Writing begins as rare, mysterious ritual technology monopolized by priests and scribes, communicating with gods and declaring kingship; bookkeeping is a later sideline. This explains early writing's remoteness and low literacy. Its methodological weakness is \textbf{preservation bias}. What survives four thousand years? Temples, tombs, palaces, tablets hardened in fires. Ordinary accounts on wood, bark, or leather have rotted. What we excavate is already the most solemn, sacred, expensively recorded subset. Treating ``what survived'' as ``what mattered most then'' is a common ancient-history trap. Perhaps countless mundane records existed but their evidence did not survive.

\textbf{Explanation three: absence (the trust account).}

This account changes angle: writing solves ``the person is not here.'' Face to face, you can ask follow-up questions. But a merchant sells goods three hundred kilometers away, a king sends orders to a frontier, creditors preserve an agreement until next year: speech fails because it cannot leave the body. Writing detaches language from the speaker. Strangers gain verifiable trust, contracts act across time and space, and empires, too large to govern by shouting, become possible. The account connects writing with contemporary growth in cities, trade, and states. Its weakness is reversing cause and effect: did empires need writing, or did writing make empires possible? Probably both, in a chicken-and-egg spiral. This explanation tells only one direction.

Each account holds some truth. Our methodological reminder remains: when an explanation claims everything, it has often discarded evidence in advance.

\section[Writing Reshapes the Mind]{Further Challenge: Writing Reshapes the Mind}
Now introduce the chapter's weightiest theory. In the twentieth century, scholars including Walter Ong, described here as Canadian and author of Orality and Literacy, proposed a bold claim: \textbf{writing does not merely put speech in a container; it changes how users think}. They called this field media studies. Its central idea fits one sentence: you think you are using writing, but writing is using you too.

The argument runs roughly as follows. In a purely oral culture, knowledge survives only by being remembered. Memory has preferences: rhythm, repetition, imagery. Oral traditions everywhere therefore favor formulas, rhyme, parallelism, and stories. Thought stays close to lived situations; knowledge grows on people and events. Once writing moves it outside memory, people can list, tabulate, define, and classify abstractly in ways oral cultures cannot. Mesopotamian word lists are evidence: scribes arranged the world's contents by category, trees in one column, offices in another, gods in another. It looks mundane but represents an earthquake in thought. Reality is no longer a concrete memory like ``Grandpa buried something under this tree,'' but an abstract object to divide, classify, and rank. On this theory, history, logic, and science are deeply writing's by-products: fixed words can be revisited, compared, and checked for contradiction.

The theory is fascinating, but beware two traps.

First, treating oral cultures as an inferior version. This misjudgment must be stopped, and a counterexample is ready to hand. The Homeric epics, Greek literature's founding works and this theory's contrast to literate thought, emerged from oral composition. Minds working without writing wove extraordinarily complex structures. Indigenous Australian songlines record water sources, landmarks, and stories along routes hundreds or thousands of kilometers long, precisely enough to guide desert crossings on foot. Oral cultures possess complex thought, located in memory, performance, and relationships rather than paper. The difference is technological fit, not intellectual rank.

Second is technological determinism: the belief that writing automatically brings reason, science, and democracy. It does not. Literate empires still practice propaganda and censorship. Tablets record blacklists as readily as accounts. Moreover, for millennia after writing's invention, literacy remained a tiny minority's privilege. The theoretical benefits accrued for ages to a small circle. This brings the truly history-changing question: writing's power is released when it becomes \textbf{cheap and widespread}.

Falling material costs brought that moment step by step. Follow the chain. Clay is free but heavy; one tablet cannot hold a letter. Egyptian papyrus is light, but geographically limited and expensive. Medieval Europe uses parchment: a substantial book consumes over a hundred sheepskins, making books valuable enough to chain to desks. Around 105 CE, Cai Lun improves Chinese papermaking; paper later reaches Europe through the Arab world, loosening prices. Then printing: Tang China has woodblock printing; the Diamond Sutra printed in 868 is the earliest surviving explicitly dated printed work. Bi Sheng invents movable type in the eleventh century; around 1450 Gutenberg brings its widespread use to Europe. One honest qualification: movable type long struggled against woodblocks in Chinese. Thousands of characters made casting, finding, and setting type extremely costly; a carved block could be printed repeatedly. Technology has no universally ``advanced'' form, only solutions fitted to particular languages and conditions, a principle already encountered here.

Every reduction in media cost widened literacy and shook power's foundations. Within decades of printing spreading in Europe, the Bible was translated into local languages and printed in affordable booklets. The church's monopoly over interpreting God's word began to crack; the Reformation followed. Korea offers a more direct example. Sejong created Hangul in 1443 and promulgated the Hunminjeongeum in 1446. Its preface explains his motive: twenty-eight new letters so everyone could learn easily and use them daily. In plain terms: ordinary people unable to read Chinese characters should easily write their own language. It was indeed easy to learn, yet the scholar-official elite resisted for centuries, disparaging it as a vernacular script for women and children. They understood the stakes: if everyone could write, writing's badge of privilege would break. Remember this picture: a script's inventor wants the widest possible learning, while old elites desperately resist. Literacy has always been political, never merely technical.

\section{From Typewriters to Emoji}
The chain of cheaper media and rearranged power reaches your hands.

English typewriters appeared in the late nineteenth century; their QWERTY layout remains. For Chinese, typewriters posed a nightmare: how could a twenty-six-key mechanism hold thousands of characters? Engineers and intellectuals struggled. In the 1940s, Lin Yutang spent his fortune developing the MingKwai typewriter, using ingenious shape-based retrieval to solve input, but high costs prevented mass production. The eventual solution was coding, not mechanics. Pinyin input retrieves characters by sound; Wubi decomposes and encodes their components. Every time you type, you answer section five's choice: the sound route or the shape route?

Digitization does something more fundamental: gives the world's characters identity numbers. Unicode assigns a unique number to nearly every known human character, bringing Chinese, kana, Hangul, Arabic, Devanagari, and Cherokee signs into one huge table so they can share a screen. Its newest, liveliest residents are emoji: yellow faces, flags, flames. Consider the winding path: pictures become signs, signs writing, writing numbers, and digital-age people return to chatting with little pictures. After four millennia, pictures come back as stickers. Of course, emoji are far from cuneiform: they have no fixed reading and record no single language, so our section-five tool classifies them as symbols, not writing. But they remind us that reliance on pictures never vanished.

Then literacy. Two centuries ago, readers and writers were a minority worldwide. Today adult literacy is generally recognized as above eighty percent, humanity's largest redistribution of this power, giving countless people their first right to put things in black and white. Yet new boundaries grow where old ones stood. The warehouse farmer's predicament has changed form, not disappeared. The question is no longer merely whether you can write, but whose content algorithms promote, whose speech is hidden, throttled, or deleted; whose data is minutely recorded, who is blank in network memory; on whose servers your digital footprints reside, and who can alter, sell, or erase them. More people can write, yet deciding what remains is concentrating in a few platforms. Four thousand years later, memory and power wear new clothes around the same entanglement.

A gentle irony remains. Thamus's feared forgetfulness returns in new forms. Input methods write for us, so people forget how to form familiar characters by hand: confronted with \mbox{喷嚏}, pēntì, ``sneeze,'' they cannot recall how to write \mbox{嚏}. Contacts remember numbers, so we cannot recite even our home phone. Clouds store photographs, so we less often truly recall a journey. More counterintuitively, digital writing may be more fragile than clay. Fire hardens tablets, helping them survive four thousand years; one failed drive, closed platform, or dead link can erase years of records overnight. Our era produces more writing than any other and may also produce history's most readily evaporated memory.

\section{Over to You: Outsourcing All Memory}
Return to the complaint tablet.

Nanni probably meant only to vent, perhaps intimidate a copper merchant, 3,750 years ago. He could not imagine his tablet outliving his city, kingdom, and language long enough for a teenager today to encounter his anger through glass. Writing gave that anger passage through time, a chance then offered to only a few who could write and chose a durable material.

You belong to history's first generation figuratively ``born literate'': from birth, input methods, cloud storage, search engines, and AI package humanity's writing capabilities into your hands. Memory is outsourced to machines on an unprecedented scale. So the final question is yours:

\textbf{When everything can be outsourced---contacts remember people, search remembers knowledge, clouds remember experiences, AI even writes for you---what do you still want to remember yourself? Is remembering personally still worthwhile?}

Give your answer and reasons. Before concluding, weigh the considerations again.

You have Thamus's case: what is forgotten may be part of you, not merely information; dependence on external signs may offer wisdom's appearance alone.

You have the scribe's case too: tools free brains for higher work. Knowing how to search, use, and question may matter more than memorizing. Every generation since writing has defended new tools this way.

There is also power's dimension: whoever owns the server holding memory can delete it. Could someone entrusting all memory to others wake one morning like the warehouse farmer, speechless before ``That is what the system records''?

There is no standard answer. But notice: the characters expressing your answer are being carried by some medium, stored by some system, permitted by some power. From tablet to screen, materials change; two tangled questions remain. Who remembers for you, and who decides?
